% ---
% title: PIN codes
% tags: finance, counting, permutations
% ---

\begin{problem}
	\begin{context}
		A bank card is protected by a 4-digit PIN; after three wrong attempts the
		card is blocked. The number of possible PINs measures how hard it is to
		guess one.
	\end{context}
	A bank card PIN has 4 digits. (a) How many PINs are there? (b) How many
	have no repeated digit? (c) A thief makes three different random guesses.
	Which fraction of all PINs does he cover?
\end{problem}

\begin{slides}
	\small
	\exstatement
	\pause
	\sol (a) $10^4=10\,000$. \ (b)
	\[ 10\cdot9\cdot8\cdot7=\frac{10!}{6!}=5040. \]
	(c) $3/10\,000$ --- next week this is a probability.
\end{slides}

\begin{solution}
	\idea A PIN is an ordered string of 4 digits (1234 and 4321 are different
	PINs), so order matters.

	\step{(a)} Each of the 4 places has 10 choices: $10^4=10\,000$ (from 0000 to
	9999).

	\step{(b)} Each digit used removes one option for the next place:
	\[ 10\cdot9\cdot8\cdot7=\frac{10!}{6!}=5040. \]

	\step{(c)} Three different guesses are 3 of the $10\,000$ equally possible
	PINs, so they cover a fraction $3/10\,000 = 0.0003$. Next week we will call
	this a probability.

	\begin{answer}
		(a) $10\,000$ \qquad (b) $5040$ \qquad (c) $3/10\,000=0.03\%$
	\end{answer}
\end{solution}
